% §2 Exercises
%
% Cross-references and their resolution status:
%
% Already resolved (within chapter 2, drafted in Tasks 2-5 of Plan 2):
%   \Cref{sec:bsc}                  -- §2.1 Binary Symmetric Channel
%   \Cref{def:bsc}                  -- §2.1 BSC formal definition
%   \Cref{sec:asymmetric-channel}   -- §2.2 Asymmetric Binary Channel
%   \Cref{def:asymmetric-channel}   -- §2.2 asymmetric channel definition
%   \Cref{def:z-channel}            -- §2.2 Z-channel definition
%   \Cref{sec:channel-composition}  -- §2.3 Composition of Channels
%   \Cref{def:channel-composition}  -- §2.3 serial composition definition
%   \Cref{eq:bsc-composition}       -- §2.3 BSC composition formula
%   \Cref{sec:model-order}          -- §2.4 Model Order
%   \Cref{def:model-order}          -- §2.4 model order definition
%   \Cref{thm:boolean-ceiling}      -- §2.4 Boolean ceiling theorem
%
% Already resolved (chapter 4, drafted in Plan 1):
%   \Cref{sec:axiom-2}             -- chapter 4 §4.3
%   \Cref{thm:kronecker}           -- chapter 4 §4.4
%   \Cref{cor:parameter-reduction} -- chapter 4 §4.4
%
% Unresolved forward refs (will resolve in later Plans):
%   \Cref{sec:composition-theorem} -- chapter 3 §3.3 (will resolve in Plan 3)
%   \Cref{ch:approximate-functions}-- chapter 9 Part III (will resolve in a later plan)
%
% label \label{ex:bsc-capacity} on Exercise 4 resolves the forward
% reference placed at bsc.tex line 304 (Exercise~\ref{ex:bsc-capacity}).
%
% Difficulty-marker macros (not in alex.sty; defined locally if not yet defined):
\providecommand{\exstar}{$\star$}
\providecommand{\exstarstar}{$\star\star$}
\providecommand{\exstarstarstar}{$\star\star\star$}

\section*{Exercises}
\addcontentsline{toc}{section}{Exercises}

\begin{enumerate}

% -----------------------------------------------------------------------
% WARM-UP  (★ × 3)
% -----------------------------------------------------------------------

\item \exstar
\textbf{Writing the channel matrix.}
Let $\fprate = 0.07$ and $\fnrate = 0.12$.
Write the $2 \times 2$ channel matrix $Q$ for the asymmetric binary
channel (\Cref{def:asymmetric-channel}) with these parameters.
Verify that both rows sum to 1 (row-stochastic condition).
Then write the matrix for the special case $\fprate = \fnrate = 0.07$
and identify which named channel (\Cref{def:bsc}) it gives.

\item \exstar
\textbf{Composing two BSCs by hand.}
Let $Q_1$ be a BSC with crossover $\fprate_1 = 0.05$ and $Q_2$ be a
BSC with crossover $\fprate_2 = 0.10$ (\Cref{def:bsc}).
\begin{enumerate}
  \item[(a)] Write out $Q_1$ and $Q_2$ explicitly as $2 \times 2$ matrices.
  \item[(b)] Compute the matrix product $Q_{\mathrm{total}} = Q_2\, Q_1$
    entry by entry (\Cref{def:channel-composition}).
  \item[(c)] Using the formula $\fprate_{\mathrm{total}} = \fprate_1(1-\fprate_2) +
    (1-\fprate_1)\fprate_2$ from \Cref{eq:bsc-composition}, verify that
    $\fprate_{\mathrm{total}} = 0.14$ and that $Q_{\mathrm{total}}$ is a
    BSC with this crossover.
\end{enumerate}

\item \exstar
\textbf{Identifying model order.}
For each of the following systems, state the model order
(\Cref{def:model-order}) and justify your answer in one sentence.
\begin{enumerate}
  \item[(a)] A BSC with crossover $\fprate = 0.03$ (\Cref{def:bsc}).
  \item[(b)] A Bloom filter (\Cref{def:z-channel}).
  \item[(c)] An asymmetric binary channel with $\fprate = 0.05$ and
    $\fnrate = 0.08$ (\Cref{def:asymmetric-channel}).
  \item[(d)] The composition of a BSC followed by a Z-channel
    (\Cref{def:channel-composition}).
\end{enumerate}

% -----------------------------------------------------------------------
% MEDIUM  (★★ × 5)
% -----------------------------------------------------------------------

\item \exstarstar
\label{ex:bsc-capacity}
\textbf{Deriving BSC capacity.}
The capacity of a channel is the maximum mutual information $I(X;Y)$
between input $X$ and output $Y$ over all input distributions.
Let $p = \Pr\{X = \top\}$ be the prior probability of a positive input,
so $\Pr\{X = \bot\} = 1 - p$.
\begin{enumerate}
  \item[(a)] For a BSC with crossover $\fprate$ (\Cref{def:bsc}), write out
    $\Pr\{Y = \top\}$ and $\Pr\{Y = \bot\}$ in terms of $p$ and $\fprate$.
  \item[(b)] Recall that mutual information satisfies
    $I(X;Y) = H(Y) - H(Y \Given X)$, where $H$ denotes binary entropy.
    Show that $H(Y \Given X) = H_2(\fprate)$ regardless of $p$, where
    $H_2(\fprate) = -\fprate \log_2 \fprate - (1-\fprate)\log_2(1-\fprate)$.
  \item[(c)] Show that $H(Y)$ is maximized at $p = 1/2$, and conclude
    that the BSC capacity is $1 - H_2(\fprate)$.
  \item[(d)] Evaluate the capacity at $\fprate = 0.05$ (the running example
    from \Cref{sec:bsc}) and at $\fprate = 0.5$. Interpret both values.
\end{enumerate}

\item \exstarstar
\textbf{Bloom filter as a Z-channel.}
The Bloom filter satisfies $\fnrate = 0$: it never reports a false negative
(\Cref{def:z-channel}).
\begin{enumerate}
  \item[(a)] Write the channel matrix $Q$ for a Bloom filter with false-positive
    rate $\fprate = 0.05$ and verify that the matrix satisfies $\fnrate = 0$.
  \item[(b)] A query returns ``not in set.'' Show that, under the Z-channel model
    (\Cref{sec:asymmetric-channel}), this answer carries no uncertainty: the
    true membership status is determined exactly.
  \item[(c)] Argue that no amount of serial composition of Z-channels
    (\Cref{def:channel-composition}) with $\fnrate = 0$ can produce a non-zero
    false-negative rate. Use the matrix-product form to support your argument.
\end{enumerate}
(\textit{Hint for (c):} consider what the composition formula from
\Cref{eq:bsc-composition} becomes when $\fnrate_1 = \fnrate_2 = 0$.)

\item \exstarstar
\textbf{The composition of two BSCs is a BSC.}
Prove the general case of the composition result implicit in
\Cref{eq:bsc-composition}: for any two BSC matrices $Q_1$ and $Q_2$
with crossovers $\fprate_1$ and $\fprate_2$, the product
$Q_{\mathrm{total}} = Q_2 Q_1$ is again a BSC.
\begin{enumerate}
  \item[(a)] A $2 \times 2$ stochastic matrix is a BSC (\Cref{def:bsc}) if
    and only if both of its off-diagonal entries are equal. Show that
    $Q_{\mathrm{total}}$ satisfies this condition by multiplying the two
    symmetric matrices and comparing off-diagonal entries.
  \item[(b)] Identify the crossover of $Q_{\mathrm{total}}$ in terms of
    $\fprate_1$ and $\fprate_2$.
  \item[(c)] The composition of BSCs is therefore \emph{closed} in the class
    of BSCs. Give a one-sentence interpretation: why is closure a useful
    property for building serial pipelines out of symmetric components?
\end{enumerate}

\item \exstarstar
\textbf{When does an asymmetric channel reduce to a BSC?}
An asymmetric binary channel (\Cref{def:asymmetric-channel}) has parameters
$\fprate$ and $\fnrate$.
\begin{enumerate}
  \item[(a)] State the algebraic condition on $(\fprate, \fnrate)$ that makes
    the asymmetric channel's matrix equal to a BSC matrix (\Cref{def:bsc}).
  \item[(b)] In the $(\fprate, \fnrate)$ parameter space $[0,1)^2$, describe
    the set of parameter pairs that correspond to BSCs. What geometric object
    is this set?
  \item[(c)] At the boundary of this set, perturb $\fnrate$ by a small amount
    $\delta > 0$ while fixing $\fprate$. Compute the new $Q$ and identify the
    model order (\Cref{def:model-order}) of the resulting channel.
  \item[(d)] Explain in one paragraph why the BSC is a measure-zero special
    case within the space of asymmetric channels, and what that implies for
    real-world data systems.
\end{enumerate}

\item \exstarstar
\textbf{Constructing a model-order-2 channel from two model-order-1 channels.}
Two deterministic channels each have model order 1: the identity channel
(crossover $\fprate = 0$) and the complementing channel (crossover $\fprate = 1$).
Both are trivially order-1 BSCs.
\begin{enumerate}
  \item[(a)] Compute the serial composition of the identity channel followed
    by the complementing channel (\Cref{def:channel-composition}).
    What channel does the composition produce?
  \item[(b)] Now consider a non-degenerate construction: compose a BSC with
    crossover $\fprate_1 = 0.1$ with an asymmetric channel having
    $\fprate_2 = 0.2$, $\fnrate_2 = 0.3$.
    Compute $Q_{\mathrm{total}} = Q_2 Q_1$ and verify the model order
    (\Cref{def:model-order}) of the result.
  \item[(c)] Describe in general how serial composition (\Cref{sec:channel-composition})
    affects the partition structure. Does composing two order-$n$ channels
    always yield an order-$n$ channel? Does it always yield an order-$m$
    channel for some fixed $m$ related to the input domain?
\end{enumerate}
(\textit{Hint:} focus on the input domain of the composed channel, not
the intermediate domain.)

% -----------------------------------------------------------------------
% CHALLENGING / EXTENSION  (★★★ × 3)
% -----------------------------------------------------------------------

\item \exstarstarstar
\textbf{Rigorous proof of the Boolean ceiling.}
\Cref{thm:boolean-ceiling} states that any binary channel over the input
domain $\{\top, \bot\}$ has model order at most 2. Give a self-contained
proof from the definitions.
\begin{enumerate}
  \item[(a)] State the definition of model order (\Cref{def:model-order})
    precisely: a partition of $\{\top, \bot\}$ is a collection of non-empty
    pairwise-disjoint subsets whose union is $\{\top, \bot\}$. List all
    distinct partitions of the two-element set $\{\top, \bot\}$.
  \item[(b)] Show that no partition of $\{\top, \bot\}$ can have more than
    2 blocks. (The argument is purely combinatorial: a partition into $k$
    blocks of a 2-element set requires $k \leq 2$.)
  \item[(c)] Conclude that the model order of any channel over $\{\top, \bot\}$
    is at most 2, and give the two distinct channel classes that correspond
    to model order 1 (the BSC) and model order 2 (the general asymmetric
    channel).
  \item[(d)] Contrast with the n-ary case: state, without proving, how the
    ceiling changes for an input domain of cardinality $k$ (foreshadows
    \Cref{ch:approximate-functions} without entering Part III content).
\end{enumerate}

\item \exstarstarstar
\textbf{Channel composition: associative but not commutative.}
Let $Q_1$, $Q_2$, $Q_3$ be binary channel matrices and define serial
composition as matrix product (\Cref{def:channel-composition}).
\begin{enumerate}
  \item[(a)] \textit{Associativity.} Prove that
    $(Q_3 Q_2) Q_1 = Q_3 (Q_2 Q_1)$.
    (This follows directly from associativity of matrix multiplication;
    state the argument clearly and identify what it implies for building
    pipelines of arbitrary length.)
  \item[(b)] \textit{Non-commutativity by counterexample.}
    Let $Q_1$ be a BSC with crossover $\fprate_1 = 0.1$ and $Q_2$ be an
    asymmetric channel with $\fprate_2 = 0.4$ and $\fnrate_2 = 0.1$.
    Compute $Q_2 Q_1$ and $Q_1 Q_2$ explicitly. Show that they are not equal.
  \item[(c)] Interpret the non-commutativity result: the order in which
    two noisy stages appear in a serial pipeline matters. Give a one-sentence
    practical implication for designing a two-stage classification pipeline
    where one stage has higher false-positive rate and the other has higher
    false-negative rate.
\end{enumerate}

\item \exstarstarstar
\textbf{A 3-input partition and model order 3.}
The binary domain $\{\top, \bot\}$ supports at most model order 2
(\Cref{thm:boolean-ceiling}). This exercise explores what happens when the
input domain grows.
\begin{enumerate}
  \item[(a)] Let the input domain be $\{1, 2, 3\}$ (a 3-element set). List
    all distinct partitions of this domain. How many distinct values of model
    order are possible?
  \item[(b)] Write a $3 \times 2$ channel matrix $Q$ for a channel from the
    3-element input domain to the binary output domain $\{\top, \bot\}$,
    choosing a partition into 2 blocks $B_1 = \{1, 2\}$ and $B_2 = \{3\}$
    with per-block error rates $\fprate_1 = 0.05$, $\fprate_2 = 0.20$.
    Verify that both rows in each block share the same entries and that the
    matrix is row-stochastic.
  \item[(c)] Write a second channel matrix using the 3-block partition
    $B_1 = \{1\}$, $B_2 = \{2\}$, $B_3 = \{3\}$ with rates
    $\fprate_1 = 0.05$, $\fprate_2 = 0.15$, $\fprate_3 = 0.30$.
    State the model order of this channel.
  \item[(d)] Explain in one paragraph why the Kronecker factorization
    (\Cref{thm:kronecker}) discussed in the context of parallel queries
    (chapter 4 §4.4) extends naturally to n-ary input domains, and identify
    the one change in the parameter count. (This is a preview of
    \Cref{ch:approximate-functions}; no detailed proof is expected.)
\end{enumerate}

\end{enumerate}
